%% =====================================================================
%%  T-14-02  The Contrast Principle
%%  -------------------------------------------------------------------
%%  One idea per file. Core is mandatory. Slots in canonical order.
%%  Max 700 words. No layout commands. No filler. New source material
%%  for this entry goes in sources/<BOOK>/T-14-02.tex -- never by
%%  rewriting this file (docs/RULES.md R0.1).
%% =====================================================================
%% @kind: topic
%% @id: T-14-02
%% @title: The Contrast Principle
%% @subtitle: A second option is judged against the first, not on its own
%% @chapter: c14
%% @order: 20
%% @status: stable
%% @sources: B4
%% @updated: 2026-09-19
%% @allow-rewrite: slot migration to the seven-question format (docs/RULES.md section 2)

\topic{T-14-02}{The Contrast Principle}{A second option is judged against the first, not on its own}


\begin{core}
Judgements are comparative. The same price, the same request and the same concession are evaluated differently depending on what was presented immediately before, so whoever chooses the order of presentation chooses a good part of the verdict.
\end{core}

\begin{mechanism}
There is no absolute scale available for most decisions. What counts as a fair price, a large favour or a serious problem is not knowable in isolation, so the mind settles it by comparison with whatever was most recently in view. That makes the first item a measuring instrument rather than a candidate: it is frequently not meant to be accepted at all, and its function is to set the distance against which the second item looks small.

The principle is not confined to price. Severity, urgency and generosity all take their scale from the immediately preceding case, which is why an offer can be made to look reasonable without changing the offer --- only what precedes it. The control lies entirely with the party who decides the sequence.
An agent who shows two unsuitable properties before the one she means to let has not wasted an afternoon; she has built the ruler the third will be measured against.
\end{mechanism}


\begin{application}
  \item Set the anchor deliberately, and let it introduce the proposal you actually want accepted.
  \item Present the expensive case first whenever the cheap case has to be accepted on its own terms.
  \item Move in visible steps rather than arriving at your real position immediately; a position reached by stages is judged against the stages.
\end{application}

\begin{feedback}
  \item Two or more options arrive in a fixed order, and the first is described at length and then abandoned without any effort to sell it.
  \item You find the second option reasonable and cannot say what you are comparing it against.
  \item The comparison set is supplied entirely by the other party.
  \item An item you would have refused a week ago now looks modest.
\end{feedback}

\begin{failure}
Anchoring is so ordinary that using it deliberately changes little in most transactions, and refusing to be anchored at all makes you slower and more suspicious than the situation calls for. The exposure worth guarding is narrow: it is where the anchor is far from the truth and the other party controls the entire comparison set. That is why the defence is to bring a comparison in from outside rather than to argue about the one on offer --- arguing inside a frame accepts the frame.
\end{failure}

\begin{countermeasures}
  \item Ask for the options out of order, or for each one separately.
  \item Import your own comparison. Price the thing against the market, not against the alternative on the table.
  \item Evaluate the final proposal as though it were the first thing you had heard, which is what it would have been in an honest sequence.
  \item Notice the abandoned first option. An offer presented at length and then dropped without argument was probably never for sale.
\end{countermeasures}

\sources{B4}
\seesources
\seealso{T-14-01, T-15-01, T-25-01}
